\documentclass[11pt]{article}
\usepackage[hmargin=0.85in,vmargin=0.8in]{geometry}
\usepackage{amsmath,amssymb,amsthm}
\usepackage[T1]{fontenc}
\usepackage{lmodern}
\usepackage[hidelinks]{hyperref}
\setlength{\emergencystretch}{1em}
\newtheorem{theorem}{Theorem}
\newtheorem{lemma}{Lemma}
\newcommand{\F}{\mathbb F}
\newcommand{\Pow}{\mathcal P}
\newcommand{\Ccal}{\mathcal C}
\title{An Executable Exhaustive Catalogue for Conjecture 00000001029}
\author{idealistichacker \quad (AI-assisted)}
\date{October 29, 2026}
\begin{document}
\maketitle
\begin{abstract}
For every prime power $q$ and all natural $n,k$, a single terminating
algorithm constructs a finite field of order $q$ from operation tables
and returns a duplicate-free list of labelled $S(2,k,n)_q$ designs.
Membership is equivalent to the actual subspace-design condition, and
every design over any field of order $q$ occurs after a fixed field
isomorphism. This proves the source's finite-explicit-list assertion,
in particular when $k\mid n$. It is exhaustive enumeration, not an
efficient algorithm, a nonemptiness theorem, or an isomorphism classification.
\end{abstract}

\section{Statement and Meaning of the Output}
Both source languages request a finite explicit list of $S(2,k,n)_q$
for each $n$ and $k\mid n$. Here $q$ has its standard finite-field meaning:
$q=p^m$ for a prime $p$ and $m>0$. Let $K$ be a field of order $q$ and
$V=K^n$. A design is a set $\mathcal D$ of linear subspaces of $V$ such that
\begin{equation}
  2\leq k\leq n,\qquad
  \forall W\in\mathcal D,\ \dim_KW=k,\qquad
  \forall U\leq V,\ \dim_KU=2\Longrightarrow
    \exists!W\in\mathcal D,\ U\subseteq W.       \tag{D}
\end{equation}
The ambient space is finite, so every set of its subspaces is finite.
This is proved rather than assumed for arbitrary block families.

\begin{theorem}\label{thm:main}
There is an executable function $A(q,n,k)$, with no input field or
oracle. If $q$ is a prime power, it returns a field table $T$ on $q$
labels and a finite duplicate-free list $L$ of families of vector-carrier
codes. For every field $K$ of order $q$, fix a field isomorphism from
the field described by $T$ to $K$, applied in every coordinate.
Then an arbitrary set $\mathcal D$ of subspaces of $K^n$ satisfies
\textup{(D)} if and only if exactly one member of $L$ decodes and
transports to $\mathcal D$. If $2\leq k\leq n$ fails, $L$ is empty.
\end{theorem}
The theorem holds for all $n,k$, so its $k\mid n$ specialization gives
the requested assertion. No quotient by linear, semilinear or incidence
isomorphism is taken: distinct labelled designs remain distinct.
The statement contains no bound on running time or output size.
The theorem supplies the specified finite computation and proves its
exact output criterion; it does not replace computation by classical
selection of an unknown list.

\section{Constructing the Field}
Enumerate $z,o\in\{0,\ldots,q-1\}$ and all pairs of binary operation
tables $a,b$ on these labels, in the finite tuple order used by
\texttt{allFunctions}. Retain the first table satisfying:
\begin{itemize}
\item $z\ne o$; addition is associative, $z+a=a$, and each element
has a left additive inverse;
\item multiplication is associative and commutative, $o a=a$, and
multiplication distributes over addition;
\item for each $a\ne z$ there exists $b$ with $ab=o$.
\end{itemize}
All tests are finite. These are sufficient minimal field axioms,
not merely necessary tests: Mathlib's \texttt{Field.ofMinimalAxioms}
constructs an actual \texttt{Field} from them. Additive negatives and
nonzero inverses are defined by the least successful bounded search;
the inverse of zero is explicitly zero. \texttt{TableElem} is a fresh
wrapper for the labels, avoiding the unrelated modular operations on
Lean's \texttt{Fin q}. Its cardinality is proved to be $q$.

\begin{lemma}\label{lem:field}
For every prime power $q$, this table search succeeds. It uses only $q$,
and its returned operations define a computable field of cardinality $q$.
\end{lemma}
\begin{proof}
Write $q=p^m$ with $p$ prime and $m>0$. Mathlib constructs a finite
field $\operatorname{GF}(p,m)$ with exactly $p^m$ elements. Relabel its
elements by $\{0,\ldots,q-1\}$. The transported zero, one, addition
and multiplication form one of the enumerated tables and pass every
test. Thus a first successful table exists. The search and inverse
searches are executable finite operations. The possibly noncomputable
Galois-field construction and its relabelling occur only in this
existence proof, not in the algorithm's dependency path.
\end{proof}
For arbitrary natural $q$ the algorithm is total: it returns failure
if no valid table is found. No field-existence assertion is made for
$q=0$ or $q=1$.

\section{Actual Subspaces and the Incidence Filter}
Work now over the selected field $K_0$ of order $q$. Enumerate the
finite vector set $X=K_0^n$. Among all subsets $S\subseteq X$, retain
those containing zero and closed under addition and multiplication
by every scalar. Such an $S$ decodes to the actual submodule
$\{x\in V:x\in S\}$. Conversely, the carrier of every actual
submodule is one of these codes.

\begin{lemma}\label{lem:codes}
Encoding and decoding give inverse bijections between closed carrier
codes and actual subspaces. They preserve equality and containment.
For a closed carrier code $S$ and its subspace $W$,
\[ |S|=|W|=q^{\dim W}. \]
Consequently, $|S|=q^d$ if and only if $\dim W=d$.
\end{lemma}
\begin{proof}
The zero, addition and scalar closure tests are exactly the submodule
axioms, so decoding gives a submodule and encoding returns its original
carrier. The two inverse identities imply equality and containment
correspondence. The carrier and the subtype representing its subspace
are in bijection. The finite-vector-space cardinality theorem gives
$|W|=q^{\dim W}$. Since every field has $q\geq2$, the function
$d\mapsto q^d$ is injective on natural numbers.
\end{proof}

Let $E_d$ consist of the closed carrier codes of cardinality $q^d$.
When $2\leq k\leq n$, enumerate all $B\subseteq E_k$ and retain
precisely those satisfying
\begin{equation}
 \forall U\in E_2,\quad
  \bigl|\{W\in B:U\subseteq W\}\bigr|=1.       \tag{I}
\end{equation}
Outside this parameter guard return the empty list. The guard occurs
in both the running algorithm and the actual target \textup{(D)}.
There is no exceptional vacuous design for $(n,k)=(0,0)$.

\begin{lemma}\label{lem:catalogue}
Every retained $B$ decodes to a design satisfying \textup{(D)}.
Every actual design has exactly one retained labelled code family.
\end{lemma}
\begin{proof}
By Lemma~\ref{lem:codes}, $E_k$ is exactly the set of $k$-subspaces
in carrier form, and $E_2$ is exactly the set of $2$-subspaces in
that form. Containment of codes is actual subspace containment.
Cardinality one in \textup{(I)} is equivalent to existence and
uniqueness of the containing block, proving soundness.
For completeness, encode each block of an actual design; the same
equivalences establish \textup{(I)}. Applying the two inverse
encoding identities to finite families proves uniqueness. For an
arbitrary set of blocks, inject its subspaces into the finite set of
vector subsets; this proves finiteness and justifies the finite-family
representation without an extra hypothesis.
\end{proof}

\section{Executable Lists and Change of Field}
The returned list is not formed by a noncomputable selection from a
finite set. Enumerate coordinate tuples with \texttt{allFunctions}.
Enumerate their sublists, convert each to a finite carrier set, remove
duplicates, and apply the closure and cardinality tests to obtain a
list of $E_k$. Enumerate its sublists in the same way to get all
block families, filter by \textup{(I)}, and remove duplicates.
For any finite set $S$, it occurs among the sublist-to-set codes of
a list $x$ if and only if all elements of $S$ occur in $x$: the
reverse direction uses the sublist obtained by filtering $x$ by $S$.
This proves that the executable list has exactly the finite-set
catalogue's membership relation. Explicit duplicate removal proves
the returned list is duplicate-free.

Now let $K$ be any field with $q$ elements. Finite fields of the same
cardinality are isomorphic; Mathlib's
\texttt{FiniteField.ringEquivOfCardEq} supplies such an isomorphism
$e:K_0\simeq K$ for the semantic proof. Applied coordinatewise, it is
a semilinear equivalence $K_0^n\simeq K^n$. Its induced subspace
order isomorphism preserves containment and carrier cardinality.
Lemma~\ref{lem:codes} and injectivity of powers of $q$ show that it
also preserves dimension. Transporting finite block families is a
bijection, hence preserves the unique-containing-block condition in
both directions. Together with Lemma~\ref{lem:catalogue}, this proves
Theorem~\ref{thm:main} for every field model. Neither $K$ nor $e$ is
an input to the running catalogue.

\section{Formalization and Verification Scope}
The Lean source uses Lean 4.33.1 and Mathlib commit
\texttt{0df444a360eaa60ab8c11dca51a86af692955474}.
The main all-model conclusion is
\texttt{QOnly.source\_conjecture}; arbitrary set-valued block families
are covered by \texttt{SetFamilies.source\_conjecture\_sets}.
\texttt{QOnly.catalogue} is the executable entry point. Source
definitions, dimension/cardinality bridges, field construction,
guarded incidence filtering, explicit lists, and model transport are
separate modules. Every advertised theorem is accompanied by actual
type and axiom output. Only the standard logical axioms
\texttt{propext}, \texttt{Classical.choice}, and \texttt{Quot.sound}
are permitted; no extra axioms or proof placeholders are used.

Strict local evaluation of the actual $q$-only entry returns one
family containing one block for $(q,n,k)=(2,2,2)$, the expected whole
space design. The invalid input $(2,0,0)$ returns an empty catalogue;
$q=0,1$ return no field model. These are small runtime diagnostics,
not substitutes for the general proofs. A prior combined diagnostic
including $n=3$ exceeded its 180-second budget; no successful $n=3$
evaluation is claimed, and that failure is retained. No large
parameter catalogue is claimed to have been computed.

The source is conjecture \texttt{00000001029}, blob
\texttt{001aa6f79c0aaaa70d0907eb42b9fb5169beb631}, checked in both
languages. This manuscript establishes a literal exhaustive-list
result. It makes no claim of structural classification, a list
with one representative per isomorphism class, efficient search,
nonexistence or existence of a design at any other parameter, first
solve credit, or upstream acceptance.

\begin{thebibliography}{9}
\raggedright
\bibitem{mathlib}
The Mathlib community, \emph{Mathlib}, fixed revision cited above.
Used results include \texttt{Field.ofMinimalAxioms},
\texttt{Module.card\_eq\_pow\_finrank},
\texttt{GaloisField.card}, and
\texttt{FiniteField.ringEquivOfCardEq}. All are imported from the
fixed source tree; the submitted source records the exact interfaces.
\end{thebibliography}
\end{document}
